\section{Does executive function moderate CBT outcome?}\label{sec:moderators}
This is the question on which the argument stands or falls. The direct evidence is thin and mostly from older adults.
\begin{itemize}
\item \textbf{Supports (at the level of the skill).} In older adults, poorer cognitive flexibility predicted poorer cognitive-restructuring quality and less distress reduction from restructuring, but \emph{not} overall CBT outcome: those with poorer flexibility ``can still benefit from standardised CBT'' \citep{johnco2014flex}.
\item \textbf{Supports (pilot).} In a pilot of CBT for GAD in older adults, participants whose executive scores remained low did not respond, while those whose scores improved responded as well as an intact-executive group \citep{mohlman2005ef}.
\item \textbf{Null or mixed.} No baseline executive effect was found in adult OCD \citep{vandborg2016memory}, child anxiety \citep{godovich2020executive} or brief psychotherapy for late-life depression \citep{schurr2025executive}; a systematic review of late-life anxiety found no replicated executive predictor \citep{kyriacou2023predictors}. In a meta-analysis of antidepressant response in late-life depression, only planning and organisation predicted outcome \citep{pimontel2016executive}.
\item \textbf{Opposite direction.} Better executive performance went with non-response to exposure-based CBT in paediatric OCD \citep{hybel2017executive}, and worse set-shifting went with better CBT response in late-life depression \citep{goodkind2016executive}.
\item \textbf{Adaptation helps.} Problem-solving therapy beat supportive therapy in older adults with executive impairment, with a larger advantage at greater impairment \citep{arean2010problem}, which points to adapting therapy rather than abandoning it.
\end{itemize}
Taken together, ``low executive function means CBT fails'' is not supported. A narrower claim is: executive capacity may limit uptake of the \emph{cognitive skill} without limiting overall response, because CBT also works through routes that do not need it. We found no adequately powered study of working-age adults with ADHD and anxiety.
